\documentclass{homework}

% The following commands set up the material that appears
% in the header.
\doclabel{Math F404: Homework 1}
\docauthor{Christopher}
\docdate{Due: September 2, 2026}

% Some common macros
\newcommand{\Reals}{\ensuremath{\mathbb R}}% Gives you a shortcut for writing the blackboard R for the real numbers - \RR
\newcommand{\Nats}{\ensuremath{\mathbb N}} % Gives you a shortcut for writing the blackboard N for the natural numbers - \NN
\newcommand{\Ints}{\ensuremath{\mathbb Z}} % Gives you a shortcut for writing the blackboard Z for the integer numbers - \ZZ

\begin{document}
\begin{problems}

\problem Sutherland 2.1
Suppose that $C, D$ are subsets of a set $X$. Prove that
\[
(X \setminus C) \cap D = D \setminus C.
\]
\solution
  \begin{proof}
$ $\newline In order to prove equality among sets, we must show that they are subsets of each other. 
    \textbf{Forward Case:} Let $x \in [(X \setminus C) \cap D]$. Unfolding the definitions:
    \begin{align*}
      x \in [( X \setminus C) \cap D] & \implies x \in (X \setminus C) \land x \in D \\
      & \implies (x \in X \land x \notin C) \land x \in D \\
      & \implies  x \in X \land x \notin C \land x \in D \\
      & \implies x \in D \land x \notin C && \text{$C \subseteq X$ and $D \subseteq X$}\\
      & \implies x \in [D \setminus C]
    \end{align*}
    Thus $(X \setminus C) \cap D \subseteq D \setminus C$. \newline
    \textbf{Backwards case}: Let $x \in [D \setminus C]$. Unfolding the definitions:
      \begin{align*}
        x \in [D \setminus C] &\implies x \in D \land x \notin C \\
        &\implies x \in D \land x \in X \land x \notin C && \text{$D \subseteq X \implies x \in D \land x \in X$}\\
        &\implies (x \in X \land x \notin C ) \land x \in D \\
        &\implies (X \setminus C) \cap D
      \end{align*}
      Thus $D \setminus X \subseteq (x \setminus C) \cap D$.
      From these two cases it follows that $(X \setminus C) \cap D = D \setminus C$.
\end{proof}

\problem Sutherland 2.2
Suppose that $A, V$ are subsets of a set $X$. Prove that
\[
A \setminus (V \cap A) = A \cap (X \setminus V).
\]
\solution

\problem Sutherland 2.5 (AI-free)
Suppose that $U_1, U_2$ are subsets of a set $X$ and that $V_1, V_2$ are
subsets of a set $Y$. Prove that
\[
(U_1 \times V_1) \cap (U_2 \times V_2) = (U_1 \cap U_2) \times (V_1 \cap V_2).
\]
\solution

\problem Sutherland 2.7
Suppose $\sim$ is an equivalence relation on a set $X$. Show that the equivalence classes partition
$X$ in the following sense. Let $\mathcal{C}$ denote the collection of equivalence classes. Then
$\bigcup_{C \in \mathcal{C}} C = X$, and if $C, C' \in \mathcal{C}$ then either $C = C'$ or
$C \cap C' = \emptyset$.
\solution

\problem Give examples of collections of open intervals $I_i = (a_i, b_i)$ with $i \in \Nats$ such that
\begin{subproblems}
\item $\bigcap_{i \in \Nats} I_i = [0,1]$
\item $\bigcap_{i \in \Nats} I_i = (0,1)$
\item $\bigcap_{i \in \Nats} I_i = (0,1]$
\end{subproblems}
For part (a) give a formal proof. For the remainder, just give an example and a brief
justification.
\subsol

\subsol

\subsol

\problem Sutherland 3.2 (No formal proof needed)
Let $f : \Reals \to \Reals$ be defined by $f(x) = \sin x$. Describe the sets:
\[
f([0, \pi/2]), \quad f([0,\infty)), \quad f^{-1}([0,1]), \quad f^{-1}([0,1/2]), \quad f^{-1}([-1,1]).
\]
\solution

\problem Sutherland 3.3
Suppose that $f : X \to Y$ and $g : Y \to Z$ are maps and $U \subseteq Z$.
Prove that
\[
(g \circ f)^{-1}(U) = f^{-1}(g^{-1}(U)).
\]
\solution

\problem Sutherland 3.4
Let $f : \Reals \to \Reals^2$ be defined by $f(x) = (x, 2x)$. Describe the sets:
\[
f([0,1]), \quad f^{-1}([0,1] \times [0,1]), \quad f^{-1}(D)
\]
where $D = \{(x,y) \in \Reals^2 : x^2 + y^2 \le 1\}$.
\solution

\problem Sutherland 3.8 (Make your proof as short as possible)
Let $f : X \to Y$ be a map and let $A, B$ be subsets of $X$. Prove
that $f(A \setminus B) = f(A) \setminus f(B)$ if and only if $f(A \setminus B) \cap f(B) = \emptyset$.
Deduce that if $f$ is injective then $f(A \setminus B) = f(A) \setminus f(B)$.
\solution

\end{problems}
\end{document}
